\documentclass[12pt,a4paper]{ctexbook}
\usepackage{geometry}
\geometry{margin=2.2cm}
\usepackage{setspace}
\setstretch{1.35}
\usepackage{hyperref}
\title{全唐诗选·poet.tang.4000}
\author{古典文献整理}
\date{}
\begin{document}
\maketitle
\tableofcontents
\newpage
\section{晦日宴高氏林亭}
\textbf{陳嘉言}\par\vspace{0.5em}
公子申敬愛，攜朋翫物華。\par
人是平陽客，地即石崇家。\par
水文生舊浦，風色滿新花。\par
日暮連歸騎，長川照晚霞。\par

\section{晦日重宴}
\textbf{陳嘉言}\par\vspace{0.5em}
高門引冠蓋，下客抱支離。\par
綺席珍羞滿，文場翰藻摛。\par
蓂華彫上月，柳色藹春池。\par
日斜歸戚里，連騎勒金羈。\par

\section{上元夜效小庾體}
\textbf{陳嘉言}\par\vspace{0.5em}
今夜可憐春，河橋多麗人。\par
寶馬金爲絡，香車玉作輪。\par
連手窺潘掾，分頭看洛神。\par
重城自不掩，出向小平津。\par

\section{晦日宴高氏林亭}
\textbf{周彥暉}\par\vspace{0.5em}
砌蓂收晦魄，津柳競年華。\par
既狎忘筌友，方淹投轄車。\par
綺筵迴舞雪，瓊醑泛流霞。\par
雲低上天晚，絲雨帶風斜。\par

\section{晦日重宴}
\textbf{周彥暉}\par\vspace{0.5em}
春華歸柳樹，俯景落蓂枝。\par
置驛銅街右，開筵玉浦陲。\par
林煙含障密，竹雨帶珠危。\par
興闌巾倒戴，山公下習池。\par

\section{晦日宴高氏林亭}
\textbf{高嶠}\par\vspace{0.5em}
飛觀寫春望，開宴坐汀沙。\par
積溜含苔色，晴空蕩日華。\par
歌入平陽第，舞對石崇家。\par
莫慮能騎馬，投轄自停車。\par

\section{晦日重宴}
\textbf{高嶠}\par\vspace{0.5em}
駕言尋鳳侶，乘歡俯雁池。\par
班荆逢舊識，斟桂喜深知。\par
紫蘭方出徑，黃鶯未囀枝。\par
別有陶春日，青天雲霧披。\par

\section{晦日宴高氏林亭}
\textbf{劉友賢}\par\vspace{0.5em}
春來日漸賒，琴酒逐年華。\par
欲向文通逕，先遊武子家。\par
池碧新流滿，巖紅落照斜。\par
興闌情未盡，步步惜風花。\par

\section{晦日宴高氏林亭}
\textbf{周思鈞}\par\vspace{0.5em}
早春驚柳穟，初晦掩蓂華。\par
騎出平陽里，筵開衞尉家。\par
竹影含雲密，池紋帶雨斜。\par
重惜林亭晚，上路滿煙霞。\par

\section{晦日重宴}
\textbf{周思鈞}\par\vspace{0.5em}
綺筵乘晦景，高宴下陽池。\par
濯雨梅香散，含風柳色移。\par
輕塵依扇落，流水入弦危。\par
勿顧林亭晚，方歡雲霧披。\par

\section{祭汾陰樂章}
\textbf{蘇頲}\par\vspace{0.5em}
禮物斯具，樂章乃陳。\par
誰其作主，皇考聖真。\par
對越在天，聖明佐神。\par
窅然汾上，厚澤如春。\par

\section{奉和聖製行次成皐途經先聖擒建德之所感而成詩應制}
\textbf{蘇頲}\par\vspace{0.5em}
漢東不執象，河朔方鬬龍。\par
夏滅漸寧亂，唐興終奮庸。\par
皇威正赫赫，兵氣何匈匈。\par
用武三川震，歸淳六代醲。\par
成臯覩王業，天下致人雍。\par
即此巡於岱，曾孫受命封。\par

\section{奉和聖製登蒲州逍遙樓應制}
\textbf{蘇頲}\par\vspace{0.5em}
在昔堯舜禹，遺塵成典謨。\par
聖皇東巡狩，況乃經此都。\par
樓觀紛迤邐，河山幾縈紆。\par
緬懷祖宗業，相繼文武圖。\par
尚德既無險，觀風諒有孚。\par
豈如汾水上，簫鼓事遊娛。\par

\section{奉和聖製過晉陽宮應制}
\textbf{蘇頲}\par\vspace{0.5em}
隋運與天絕，生靈厭氛昏。\par
聖期在寧亂，士馬興太原。\par
立極萬邦推，登庸四海尊。\par
慶膺神武帝，業付皇曾孫。\par
緬慕封唐道，追惟歸沛魂。\par
詔書感先義，典禮巡舊藩。\par
高殿綵雲合，春旗祥風翻。\par
率西見汾水，奔北空塞垣。\par
款曲童兒佐，依遲故老言。\par
里頒慈惠賞，家受復除恩。\par
下輦崇三教，建碑當九門。\par
孝思敦至美，億載奉開元。\par

\section{奉和姚令公溫湯舊館永懷故人盧公之作}
\textbf{蘇頲}\par\vspace{0.5em}
樹德豈孤邁，降神良竝出。\par
偉茲廊廟楨，調彼鹽梅實。\par
正悅虞垂舉，翻悲鄭僑卒。\par
同心不可忘，交臂何爲失。\par
清路荷前幸，明時稱右弼。\par
曾聯野外遊，尚記帷中密。\par
新慟情莫遣，舊遊詞更述。\par
空令還辱和，長歎知音日。\par

\section{和杜主簿春日有所思}
\textbf{蘇頲}\par\vspace{0.5em}
朝上高樓上，俯見洛陽陌。\par
搖蕩吹花風，落英紛已積。\par
美人不共此，芳好空所惜。\par
攬鏡塵網滋，當窗苔蘚碧。\par
緬懷在雲漢，良願暌枕席。\par
翻似無見時，如何久爲客。\par

\section{餞郢州李使君}
\textbf{蘇頲}\par\vspace{0.5em}
楚有章華臺，遙遙雲夢澤。\par
復聞擁符傳，及是收圖籍。\par
佳政在離人，能聲寄侯伯。\par
離懷朔風起，試望秋陰積。\par
中路悽以寒，羣山靄將夕。\par
傷心聊把袂，怊悵麒麟客。\par

\section{餞唐州高使君赴任}
\textbf{蘇頲}\par\vspace{0.5em}
永日奏文時，東風搖蕩夕。\par
浩然思樂事，翻復餞征客。\par
淮水春流清，楚山暮雲白。\par
勿言行路遠，所貴專城伯。\par

\section{曉濟膠川南入密界}
\textbf{蘇頲}\par\vspace{0.5em}
飲馬膠川上，傍膠南趣密。\par
林遙飛鳥遲，雲去晴山出。\par
落暉隱桑柘，秋原被花實。\par
慘然遊子寒，風露將蕭瑟。\par

\section{夜發三泉即事}
\textbf{蘇頲}\par\vspace{0.5em}
暗發三泉山，窮秋聽騷屑。\par
北林夜鳴雨，南望曉成雪。\par
祗詠北風涼，詎知南土熱。\par
沙溪忽沸渭，石道乍明滅。\par
宛若銀磧橫，復如瑤臺結。\par
指程賦所戀，遇虞不遑歇。\par
重纊濡莫解，懸旌凍猶揭。\par
下奔泥棧榰，上覯雲梯設。\par
[搏]頰羸馬頓，回眸惴人跌。\par
憧憧往復還，心注思逾切。\par
冉冉年將病，力困衰怠竭。\par
天彭信方隅，地勢誠斗絕。\par
忝曳尚書履，叨兼使臣節。\par
京坻有歲饒，亭障無邊孽。\par
歸奏丹墀左，騫能俟來哲。\par

\section{小園納涼即事}
\textbf{蘇頲}\par\vspace{0.5em}
煩暑避蒸鬱，居閑習高明。\par
長風自遠來，層閣有餘清。\par
散灑納涼氣，蕭條遺世情。\par
奈何誇大隱，終日繫塵纓。\par

\section{昆明池晏坐答王兵部珣三韻見示}
\textbf{蘇頲}\par\vspace{0.5em}
畫舸疾如飛，遙遙汎夕暉。\par
石黥吹浪隱，玉女步塵歸。\par
獨有銜恩處，明珠在釣磯。\par

\section{奉和聖製春臺望應制}
\textbf{蘇頲}\par\vspace{0.5em}
壯麗天之府，神明王者宅。\par
大君乘飛龍，登彼復懷昔。\par
圓闕朱光燄，橫山翠微積。\par
河汧流作表，縣聚開成陌。\par
即舊在皇家，維新具物華。\par
雲連所上居恆屬，日更時中望不斜。\par
三月滄池搖積水，萬年青樹綴新花。\par
暴嬴國此嘗圖霸，霸業後仁先以詐。\par
東破諸侯西入秦，咸陽北阪南渭津。\par
詩書焚爇散學士，高閣奢踰嬌美人。\par
事往覆輈經遠喻，春還按蹕憑高賦。\par
戎觀愛力深惟省，越厭陳方何足務。\par
清吹遙遙發帝臺，宸文耿耿照天迴。\par
伯夷位事愚臣忝，喜奏聲成鳳鳥來。\par

\section{長相思}
\textbf{蘇頲}\par\vspace{0.5em}
君不見天津橋下東流水，南望龍門北朝市。\par
楊柳青青宛地垂，桃紅李白花參差。\par
花參差，柳堪結，此時憶君心斷絕。\par

\section{蜀城哭台州樂安少府}
\textbf{蘇頲}\par\vspace{0.5em}
遠遊躋劒閣，長想屬天台。\par
萬里隔三載，此邦余重來。\par
音容曠不覩，夢寐殊悠哉。\par
邊郡饒藉藉，晚庭正回回。\par
喜傳上都封，因促傍吏開。\par
向悟海鹽客，已而梁木摧。\par
變衣寢門外，揮涕少城隈。\par
却記分明得，猶持委曲猜。\par
師儒昔訓獎，仲季時童孩。\par
服義題書篋，邀歡汎酒杯。\par
蹔令風雨散，仍迫歲時迴。\par
其道惟正直，其人信美偲。\par
白頭還作尉，黃綬固非才。\par
可歎懸蛇疾，先貽問鵩災。\par
故鄉閉窮壤，宿草生寒荄。\par
零落九原去，蹉跎四序催。\par
曩期冬贈橘，今哭夏成梅。\par
執禮誰爲賵，居常不徇財。\par
北登嵔𡻱坂，東望姑蘇臺。\par
天路本懸絕，江波復泝洄。\par
念孤心易斷，追往恨艱裁。\par
不遂卿將伯，孰云陳與雷。\par
吾衰亦如此，夫子復何哀。\par

\section{立春日侍宴內出剪綵花應制}
\textbf{蘇頲}\par\vspace{0.5em}
曉入宜春苑，穠芳吐禁中。\par
剪刀因裂素，妝粉爲開紅。\par
彩異驚流雪，香饒點便風。\par
裁成識天意，萬物與花同。\par

\section{春日芙蓉園侍宴應制}
\textbf{蘇頲}\par\vspace{0.5em}
御道紅旗出，芳園翠輦遊。\par
繞花開水殿，架竹起山樓。\par
荷芰輕薰幄，魚龍出負舟。\par
寧知穆天子，空賦白雲秋。\par

\section{奉和聖製人日清暉閣宴羣臣遇雪應制}
\textbf{蘇頲}\par\vspace{0.5em}
樓觀空煙裏，初年瑞雪過。\par
苑花齊玉樹，池水作銀河。\par
七日祥圖啓，千春御賞多。\par
輕飛傳綵勝，天上奉薰歌。\par

\section{奉和七夕宴兩儀殿應制}
\textbf{蘇頲}\par\vspace{0.5em}
靈媛乘秋發，仙裝警夜催。\par
月光窺欲渡，河色辨應來。\par
機石天文寫，針樓御賞開。\par
竊觀棲鳥至，疑向鵲橋迴。\par

\section{奉和九日幸臨渭亭登高應制得時字}
\textbf{蘇頲}\par\vspace{0.5em}
嘉會宜長日，高筵順動時。\par
曉光雲外洗，晴色雨餘滋。\par
降鶴因韶德，吹花入御詞。\par
願陪陽數節，億萬九秋期。\par

\section{游禁苑幸臨渭亭遇雪應制}
\textbf{蘇頲}\par\vspace{0.5em}
平明敞帝居，霰雪下凌虛。\par
寫月含珠綴，從風薄綺疏。\par
年驚花絮早，春夜管弦初。\par
已屬雲天外，欣承霈澤餘。\par

\section{奉和送金城公主適西蕃應制}
\textbf{蘇頲}\par\vspace{0.5em}
帝女出天津，和戎轉罽輪。\par
川經斷腸望，地與析支鄰。\par
奏曲風嘶馬，銜悲月伴人。\par
旋知偃兵革，長是漢家親。\par

\section{奉和聖製登驪山高頂寓目應制}
\textbf{蘇頲}\par\vspace{0.5em}
仙蹕御層氛，高高積翠分。\par
巖聲中谷應，天語半空聞。\par
豐樹連黃葉，函關入紫雲。\par
聖圖恢宇縣，歌賦小橫汾。\par

\section{幸白鹿觀應制}
\textbf{蘇頲}\par\vspace{0.5em}
碧虛清吹下，藹藹入仙宮。\par
松磴攀雲絕，花源接澗空。\par
受符邀羽使，傳訣注香童。\par
詎似閑居日，徒聞有順風。\par

\section{題壽安王主簿池館}
\textbf{蘇頲}\par\vspace{0.5em}
洛邑通馳道，韓郊在屬城。\par
館將花雨映，潭與竹聲清。\par
賢俊鸞棲棘，賓遊馬佩衡。\par
願言隨狎鳥，從此濯吾纓。\par

\section{扈從溫泉奉和姚令公喜雪}
\textbf{蘇頲}\par\vspace{0.5em}
清道豐人望，乘時漢主遊。\par
恩暉隨霰下，慶澤與雲浮。\par
泉暖驚銀磧，花寒愛玉樓。\par
鼎臣今有問，河伯且應留。\par

\section{秋社日崇讓園宴得新字}
\textbf{蘇頲}\par\vspace{0.5em}
鳴爵三農稔，勾龍百代神。\par
運昌叨輔弼，時泰喜黎民。\par
樹缺池光近，雲開日影新。\par
生全應有地，長願樂交親。\par

\section{奉和魏僕射秋日還鄉有懷之作}
\textbf{蘇頲}\par\vspace{0.5em}
南宮夙拜罷，東道晝遊初。\par
飲餞傾冠蓋，傳呼問里閭。\par
樹悲懸劒所，溪想釣璜餘。\par
明發輝光至，增榮駟馬車。\par

\section{武[擔]山寺}
\textbf{蘇頲}\par\vspace{0.5em}
武[擔]獨蒼然，墳山下玉泉。\par
鱉靈時共盡，龍女事同遷。\par
松柏銜哀處，幡花種福田。\par
詎知留鏡石，長與法輪圓。\par

\section{餞潞州陸長史再守汾州}
\textbf{蘇頲}\par\vspace{0.5em}
河尹政成期，爲汾昔所推。\par
不榮三入地，還美再臨時。\par
擁傳雲初合，聞鶯日正遲。\par
道傍多出餞，別有吏民思。\par

\section{餞荆州崔司馬}
\textbf{蘇頲}\par\vspace{0.5em}
茂禮雕龍昔，香名展驥初。\par
水連南海漲，星拱北辰居。\par
稍發仙人履，將題別駕輿。\par
明年徵拜入，荆玉不藏諸。\par

\section{送吏部李侍郎東歸得歸字}
\textbf{蘇頲}\par\vspace{0.5em}
陌上有光輝，披雲向洛畿。\par
賞來榮扈從，別至惜分飛。\par
泉溜含風急，山煙帶日微。\par
茂曹今去矣，人物喜東歸。\par

\section{送光祿姚卿還都}
\textbf{蘇頲}\par\vspace{0.5em}
漢室有英台，荀家寵俊才。\par
九卿朝已入，三子暮同來。\par
不授綸爲草，還司鼎用梅。\par
兩京王者宅，駟馬日應迴。\par

\section{春晚送瑕丘田少府還任因寄洛中鏡上人}
\textbf{蘇頲}\par\vspace{0.5em}
聞道還沂上，因聲寄洛濱。\par
別時花欲盡，歸處酒應春。\par
聚散同行客，悲歡屬故人。\par
少年追樂地，遙贈一霑巾。\par

\section{送賈起居奉使入洛取圖書因便拜覲}
\textbf{蘇頲}\par\vspace{0.5em}
舊國才因地，當朝史命官。\par
遺文徵闕簡，還思采芳蘭。\par
傳發關門候，觴稱邑里歡。\par
早持京副入，旋佇洛書刊。\par

\section{送常侍舒公歸覲}
\textbf{蘇頲}\par\vspace{0.5em}
朝聞講藝餘，晨省拜恩初。\par
訓胄尊庠序，榮親耀里閭。\par
朱丹華轂送，斑白綺筵舒。\par
江上春流滿，還應薦躍魚。\par

\section{興州出行}
\textbf{蘇頲}\par\vspace{0.5em}
危途曉未分，驅馬傍江濆。\par
滴滴泣花露，微微出岫雲。\par
松梢半吐月，蘿翳漸移曛。\par
旅客腸應斷，吟猿更使聞。\par

\section{邊秋薄暮}
\textbf{蘇頲}\par\vspace{0.5em}
海外秋鷹擊，霜前旅雁歸。\par
邊風思鞞鼓，落日慘旌麾。\par
浦暗漁舟入，川長獵騎稀。\par
客悲逢薄暮，況乃事戎機。\par

\section{曉發方騫驛}
\textbf{蘇頲}\par\vspace{0.5em}
傳置遠山蹊，龍鍾蹴澗泥。\par
片陰常作雨，微照已生霓。\par
鬢髪愁氛換，心情險路迷。\par
方知向蜀者，偏識子規啼。\par

\section{經三泉路作}
\textbf{蘇頲}\par\vspace{0.5em}
三月松作花，春行日漸賒。\par
竹障山鳥路，藤蔓野人家。\par
透石飛梁下，尋雲絕磴斜。\par
此中誰與樂，揮涕語年華。\par

\section{故高安大長公主挽詞}
\textbf{蘇頲}\par\vspace{0.5em}
彤管承師訓，青圭備禮容。\par
孟孫家代寵，元女國朝封。\par
柔軌題貞順，閑規賦肅雍。\par
寧知落照盡，霜吹入悲松。\par

\section{贈司徒豆盧府君挽詞}
\textbf{蘇頲}\par\vspace{0.5em}
寵贈追胡廣，親臨比賀循。\par
幾聞投劒客，多會服緦人。\par
草閉墳將古，松陰地不春。\par
二陵猶可望，存歿有忠臣。\par

\section{故右散騎常侍舒國公褚公挽詞}
\textbf{蘇頲}\par\vspace{0.5em}
陽翟疏豐構，臨平演慶源。\par
學筵尊授几，儒服寵乘軒。\par
審諭留中密，開陳與上言。\par
徂暉一不借，空有賜東園。\par

\section{奉和初春幸太平公主南莊應制}
\textbf{蘇頲}\par\vspace{0.5em}
主第山門起灞川，宸遊風景入初年。\par
鳳皇樓下交天仗，烏鵲橋頭敞御筵。\par
往往花間逢綵石，時時竹裏見紅泉。\par
今朝扈蹕平陽館，不羨乘槎雲漢邊。\par

\section{奉和春日幸望春宮應制}
\textbf{蘇頲}\par\vspace{0.5em}
東望望春春可憐，更逢晴日柳含煙。\par
宮中下見南山盡，城上平臨北斗懸。\par
細草徧承回輦處，輕花微落奉觴前。\par
宸遊對此歡無極，鳥哢聲聲入管絃。\par

\section{人日重宴大明宮恩賜綵縷人勝應制}
\textbf{蘇頲}\par\vspace{0.5em}
疏龍磴道切昭回，建鳳旗門繞帝臺。\par
七葉仙蓂依月吐，千株御柳拂煙開。\par
初年競貼宜春勝，長命先浮獻壽杯。\par
是日皇靈知竊幸，羣心就捧大明來。\par

\section{侍宴安樂公主山莊應制}
\textbf{蘇頲}\par\vspace{0.5em}
駸駸羽騎歷城池，帝女樓臺向晚披。\par
霧灑旌旗雲外出，風回巖岫雨中移。\par
當軒半落天河水，繞徑全低月樹枝。\par
簫鼓宸遊陪宴日，和鳴雙鳳喜來儀。\par

\section{興慶池侍宴應制}
\textbf{蘇頲}\par\vspace{0.5em}
降鶴池前回步輦，棲鸞樹杪出行宮。\par
山光積翠遙疑逼，水態含青近若空。\par
直視天河垂象外，俯窺京室畫圖中。\par
皇歡未使恩波極，日暮樓船更起風。\par

\section{廣達樓下夜侍酺宴應制}
\textbf{蘇頲}\par\vspace{0.5em}
東嶽封廻宴洛京，西墉通晚會公卿。\par
樓臺絕勝宜春苑，燈火還同不夜城。\par
正覩人間朝市樂，忽聞天上管弦聲。\par
酺來萬舞羣臣醉，喜戴千年聖主明。\par

\section{龍池樂章}
\textbf{蘇頲}\par\vspace{0.5em}
西京鳳邸躍龍泉，佳氣休光鎮在天。\par
軒后霧圖今已得，秦王水劒昔常傳。\par
恩魚不入昆明釣，瑞鶴長如太液仙。\par
願侍巡遊同舊里，更聞簫鼓濟樓船。\par

\section{扈從鄠杜間奉呈刑部尚書舅崔黃門馬常侍}
\textbf{蘇頲}\par\vspace{0.5em}
翠輦紅旗出帝京，長楊鄠杜昔知名。\par
雲山一一看皆美，竹樹蕭蕭畫不成。\par
羽騎將過持袂拂，香車欲度捲簾行。\par
漢家曾草巡遊賦，何似今來應聖明。\par

\section{景龍觀送裴士曹}
\textbf{蘇頲}\par\vspace{0.5em}
昔日嘗聞公主第，今時變作列仙家。\par
池傍坐客穿叢篠，樹下遊人掃落花。\par
雨雪長疑向函谷，山泉直似到流沙。\par
君還洛邑分明記，此處同來閱歲華。\par

\section{春晚紫微省直寄內}
\textbf{蘇頲}\par\vspace{0.5em}
直省清華接建章，向來無事日猶長。\par
花間燕子棲鳷鵲，竹下鵷雛繞鳳皇。\par
內史通宵承紫誥，中人落晚愛紅妝。\par
別離不慣無窮憶，莫誤卿卿學太常。\par

\section{贈彭州權別駕}
\textbf{蘇頲}\par\vspace{0.5em}
雙流脈脈錦城開，追餞年年往復廻。\par
秪道歌謠迎半刺，徒聞禮數揖中台。\par
黃鶯急囀春風盡，斑馬長嘶落景催。\par
莫愴分飛岐路別，還當奏最掖垣來。\par

\section{寒食宴于中舍別駕兄弟宅}
\textbf{蘇頲}\par\vspace{0.5em}
子推山上歌龍罷，定國門前結駟來。\par
始覩元昆鏘玉至，旋聞季子佩刀廻。\par
晴花處處因風起，御柳條條向日開。\par
自有長筵歡不極，還將綵服詠南陔。\par

\section{九月九日望蜀臺}
\textbf{蘇頲}\par\vspace{0.5em}
蜀王望蜀舊臺前，九日分明見一川。\par
北料鄉關方自此，南辭城郭復依然。\par
青松繫馬攢巖畔，黃菊留人籍道邊。\par
自昔登臨湮滅盡，獨聞忠孝兩能傳。\par

\section{奉和晦日幸昆明池應制}
\textbf{蘇頲}\par\vspace{0.5em}
炎曆事邊陲，昆明始鑿池。\par
豫遊光後聖，征戰罷前規。\par
霽色清珍宇，年芳入錦陂。\par
御杯蘭薦葉，仙仗柳交枝。\par
二石分河瀉，雙珠代月移。\par
微臣比翔泳，恩廣自無涯。\par

\section{奉和聖製幸禮部尚書竇希玠宅應制}
\textbf{蘇頲}\par\vspace{0.5em}
尚書列侯第，外戚近臣家。\par
飛棟臨青綺，迴輿轉翠華。\par
日交當戶樹，泉漾滿池花。\par
圓頂圖嵩石，方流擁魏沙。\par
豫遊今聽履，侍從昔鳴笳。\par
自有天文降，無勞訪海槎。\par

\section{奉和幸韋嗣立山莊應制}
\textbf{蘇頲}\par\vspace{0.5em}
摐金寒野霽，步玉曉山幽。\par
帝幄期松子，臣廬訪葛侯。\par
百工徵往夢，七聖扈來遊。\par
斗柄乘時轉，臺階捧日留。\par
樹重巖籟合，泉迸水光浮。\par
石徑喧朝履，璜溪擁釣舟。\par
恩如犯星夜，歡擬濟河秋。\par
不學堯年隱，空令傲許由。\par

\section{奉和聖製送張說上集賢學士賜宴得茲字}
\textbf{蘇頲}\par\vspace{0.5em}
肅肅金殿裏，招賢固在茲。\par
鏘鏘石渠內，序拜亦同時。\par
宴錫歡談道，文成貴說詩。\par
用儒今作相，敦學舊爲師。\par
下際天光近，中來帝渥滋。\par
國朝良史載，能事日論思。\par

\section{奉和聖製途經華嶽應制}
\textbf{蘇頲}\par\vspace{0.5em}
朝望蓮華嶽，神心就日來。\par
晴觀五千仞，仙掌拓山開。\par
受命金符叶，過祥玉瑞陪。\par
霧披乘鹿見，雲起馭龍廻。\par
偃樹枝封雪，殘碑石冐苔。\par
聖皇惟道契，文字勒巖隈。\par

\section{奉和聖製經河上公廟應制}
\textbf{蘇頲}\par\vspace{0.5em}
河流無日夜，河上有神仙。\par
輦路曾經此，壇場即宛然。\par
下疑成洞穴，高若在空煙。\par
善物遺方外，和光繞道邊。\par
事因周史得，言與漢王傳。\par
喜屬膺期聖，邦家業又玄。\par

\section{奉和聖製答張說出雀鼠谷}
\textbf{蘇頲}\par\vspace{0.5em}
雨施巡方罷，雲從訓俗廻。\par
密途汾水衞，清蹕晉郊陪。\par
寒著山邊盡，春當日下來。\par
御祠玄鳥應，仙仗綠楊開。\par
作頌音傳雅，觀文色動台。\par
更知西向樂，宸藻協鹽梅。\par

\section{奉和恩賜樂遊園宴應制}
\textbf{蘇頲}\par\vspace{0.5em}
樂遊光地選，酺飲慶天從。\par
座密千官盛，場開百戲容。\par
綠塍際山盡，緹幕倚雲重。\par
下上花齊發，周廻柳遍濃。\par
奪晴紛劒履，喧聽雜歌鐘。\par
日晚銜恩散，堯人倂可封。\par

\section{恩制尚書省僚宴昆明池同用堯字}
\textbf{蘇頲}\par\vspace{0.5em}
露渥灑雲霄，天官次斗杓。\par
昆明四十里，空水極晴朝。\par
雁似銜紅葉，鯨疑噴海潮。\par
翠山來徹底，白日去廻標。\par
泳廣漁杈溢，浮深妓舫搖。\par
飽恩皆醉止，合舞共歌堯。\par

\section{奉和聖製幸望春宮送[朔方]大總管張仁亶}
\textbf{蘇頲}\par\vspace{0.5em}
北風吹早雁，日夕渡河飛。\par
氣冷膠應折，霜明草正腓。\par
老臣帷幄筭，元宰廟堂機。\par
餞飲廻仙蹕，臨戎解御衣。\par
軍裝乘曉發，師律候春歸。\par
方佇勳庸盛，天詞降紫微。\par

\section{奉和聖製登太行山中言志應制}
\textbf{蘇頲}\par\vspace{0.5em}
北山東入海，馳道上連天。\par
順動三光柱，登臨萬象懸。\par
俛觀河內邑，平指洛陽川。\par
按蹕夷關險，張旗亙井泉。\par
曉巖中警柝，春事下蒐田。\par
德重周王問，歌輕漢后傳。\par
宸遊鋪令典，睿思起芳年。\par
願以封書奏，廻鑾禪肅然。\par

\section{奉和聖製漕橋東送新除岳牧}
\textbf{蘇頲}\par\vspace{0.5em}
寶賢不遺俊，臺閣盡鵷鸞。\par
未若調人切，其如簡帝難。\par
上才膺出典，中旨念分官。\par
特以專城貴，深惟列郡安。\par
政行思務本，風靡屬勝殘。\par
有令田知急，無分獄在寬。\par
至言題睿札，殊渥灑仙翰。\par
詔餞三台降，朝榮萬國歡。\par
舉杯臨水發，張樂擁橋觀。\par
式佇東封會，鏘鏘檢玉壇。\par

\section{奉和聖製途次舊居應制}
\textbf{蘇頲}\par\vspace{0.5em}
潞國臨淄邸，天王別駕輿。\par
出潛離隱際，小往大來初。\par
東陸行春典，南陽即舊居。\par
約川星罕駐，扶道日旂舒。\par
雲覆連行在，風廻助掃除。\par
木行城邑望，臯落土田疏。\par
昔試邦興后，今過俗徯予。\par
示威寧校獵，崇讓不陳魚。\par
府吏趨宸扆，鄉耆捧帝車。\par
帳傾三飲處，閑整六飛餘。\par
盛業銘汾鼎，昌期應洛書。\par
願陪歌賦末，留比蜀相如。\par

\section{奉和聖製至長春宮登樓望稼穡之作}
\textbf{蘇頲}\par\vspace{0.5em}
帝迹奚其遠，皇符之所崇。\par
敬時堯務作，盡力禹稱功。\par
赫赫惟元后，經營自左馮。\par
變蕪秔稻實，流惡水泉通。\par
國阜猶前豹，人疲詎昔熊。\par
黃圖巡沃野，清吹入離宮。\par
是閱京坻富，仍觀都邑雄。\par
憑軒一何綺，積溜寫晴空。\par
禮節家安外，和平俗在中。\par
見龍垂渭北，辭雁指河東。\par
睿思方居鎬，宸遊若飲豐。\par
寧誇子雲從，秪爲獵扶風。\par

\end{document}